\documentclass[11pt,a4paper]{article}

\usepackage[utf8]{inputenc}
\usepackage[indonesian]{babel}
\usepackage[margin=2cm]{geometry}
\usepackage{amsmath,amssymb,amsfonts}
\newenvironment{proof}{\noindent\textbf{Bukti.}}{\hfill$\blacksquare$\vspace{0.3em}}
\usepackage{booktabs}
\usepackage{longtable}
\usepackage{multirow}
\usepackage{enumitem}
\usepackage{titlesec}
\usepackage{float}
\usepackage{xcolor}
\usepackage{tcolorbox}
\usepackage{fancyhdr}
\usepackage{hyperref}

\hypersetup{
    colorlinks=true,
    linkcolor=blue,
    urlcolor=blue,
    pdftitle={12 Soal Latihan dan Pembahasan Lengkap: Teori Shannon},
    pdfauthor={MA4151 Kriptografi}
}

\pagestyle{fancy}
\fancyhf{}
\rhead{\textbf{MA4151 Kriptografi}}
\lhead{12 Soal Latihan \& Pembahasan Teori Shannon}
\rfoot{Halaman \thepage}
\setlength{\headheight}{14pt}

% Color Definitions
\definecolor{boxheader}{RGB}{30, 55, 153}
\definecolor{boxbg}{RGB}{248, 249, 250}
\definecolor{colmudah}{RGB}{38, 130, 70}
\definecolor{colsbgmudah}{RGB}{243, 250, 245}
\definecolor{colsedang}{RGB}{25, 85, 155}
\definecolor{colsbgsedang}{RGB}{242, 247, 253}
\definecolor{colsulit}{RGB}{165, 30, 30}
\definecolor{colsbgsulit}{RGB}{253, 242, 242}

\tcbset{
    arc=1.5mm,
    left=3mm,
    right=3mm,
    top=2.5mm,
    bottom=2.5mm,
    fonttitle=\bfseries
}

% Boxes for Problem Tiers
\newtcolorbox{soalmudah}[1]{
    colback=colsbgmudah,
    colframe=colmudah,
    title={#1}
}

\newtcolorbox{soalsedang}[1]{
    colback=colsbgsedang,
    colframe=colsedang,
    title={#1}
}

\newtcolorbox{soalsulit}[1]{
    colback=colsbgsulit,
    colframe=colsulit,
    title={#1}
}

\newtcolorbox{tipbox}[1]{
    colback=boxbg,
    colframe=boxheader,
    title={#1}
}

\title{\vspace{-1cm}\textbf{\Large KUMPULAN 12 SOAL LATIHAN \& PEMBAHASAN MATEMATIS LENGKAP}\\\large Topik 3: Teori Shannon (\textit{Shannon's Theory of Secrecy Systems})\\\large\textbf{Mata Kuliah MA4151 Kriptografi}}
\author{\textbf{Disusun Berdasarkan Slide Kuliah \& Catatan Komprehensif Teori Shannon}}
\date{}

\begin{document}

\maketitle
\vspace{-0.5cm}

\begin{abstract}
Dokumen ini berisi \textbf{12 soal latihan komprehensif} beserta \textbf{pembahasan matematis langkah demi langkah} untuk materi \textit{Teori Shannon} (Claude Shannon, 1949, \textit{``Communication Theory of Secrecy Systems''}). Soal-soal dikelompokkan secara proporsional ke dalam \textbf{3 tingkatan kesulitan} (masing-masing 4 soal):
\begin{enumerate}[leftmargin=*]
    \item \textbf{Tingkat Mudah (Soal 1--4)}: Entropi dasar, kompresi Huffman, jarak unicity sandi klasik sederhana, dan konsep idempotensi.
    \item \textbf{Tingkat Sedang (Soal 5--8)}: Pengujian keamanan sempurna matriks $3 \times 3$, perhitungan \textit{Key Equivocation} $H(K|C)$, jarak unicity sandi polialfabetik \& permutasi, serta komutativitas sandi affine.
    \item \textbf{Tingkat Sulit (Soal 9--12)}: Pembuktian formal syarat Shannon $|\mathcal{K}| \ge |\mathcal{C}| \ge |\mathcal{P}|$, pembuktian ketaksamaan subaditivitas entropi via Ketaksamaan Jensen, jarak unicity \& ekspektasi kunci palsu ($\bar{s}_n$) pada Sandi Hill $3 \times 3$ ($\mathrm{GL}(3, \mathbb{Z}_{26})$), serta analisis arsitektur \textit{Substitution-Permutation Network} (SPN) non-idempoten Shannon.
\end{enumerate}
\end{abstract}

\vspace{0.2cm}
\hrule
\vspace{0.3cm}

\tableofcontents

\vspace{0.4cm}
\hrule
\vspace{0.6cm}

% =========================================================================
\section{Tingkat 1: Tingkat Mudah (Soal 1 -- 4)}
% =========================================================================

\subsection{Soal 1: Entropi Variabel Acak \& Batas Entropi Maksimum}

\begin{soalmudah}{Soal 1 (Tingkat Mudah) -- Entropi Sumber Pesan \& Distribusi Seragam}
Diberikan suatu sumber informasi diskret yang memancarkan 4 kemungkinan simbol pesan $\mathcal{P} = \{A, B, C, D\}$ dengan distribusi peluang sebagai berikut:
\[
p(A) = \frac{1}{2}, \quad p(B) = \frac{1}{4}, \quad p(C) = \frac{1}{8}, \quad p(D) = \frac{1}{8}.
\]
\textbf{Pertanyaan}:
\begin{enumerate}[label=\alph*)]
    \item Hitunglah nilai entropi Shannon $H(X)$ dari variabel acak tersebut secara eksak dalam satuan bit.
    \item Tentukan nilai entropi maksimum $H_{\max}(X)$ yang dapat dicapai oleh sumber dengan 4 simbol tersebut, dan jelaskan distribusi peluang seperti apa yang menghasilkannya.
    \item Hitunglah efisiensi informasi dari distribusi tersebut yang didefinisikan sebagai $\eta = \frac{H(X)}{H_{\max}(X)}$ serta tentukan redundansi sumbernya $R = 1 - \eta$.
\end{enumerate}
\end{soalmudah}

\noindent\textbf{Pembahasan Lengkap}:

\noindent\textbf{Langkah 1: Rumus yang Digunakan}
\begin{enumerate}
    \item \textbf{Definisi Entropi Shannon}:
    \begin{equation}
    \label{eq:ent_def1}
    H(X) = - \sum_{i=1}^n p(x_i) \log_2 p(x_i) = \sum_{i=1}^n p(x_i) \log_2\left(\frac{1}{p(x_i)}\right)
    \end{equation}
    \item \textbf{Teorema Batas Maksimum Entropi}:
    \begin{equation}
    \label{eq:ent_max1}
    H_{\max}(X) = \log_2 n \quad (\text{tercapai saat } p_i = 1/n, \, \forall i)
    \end{equation}
\end{enumerate}

\noindent\textbf{Langkah 2: Perhitungan Bagian a (Entropi $H(X)$)}
Substitusikan nilai probabilitas masing-masing simbol:
\begin{itemize}
    \item Simbol $A$: $p(A) = \frac{1}{2} = 2^{-1} \implies \log_2(1/p(A)) = \log_2 2 = 1$ bit.
    \item Simbol $B$: $p(B) = \frac{1}{4} = 2^{-2} \implies \log_2(1/p(B)) = \log_2 4 = 2$ bit.
    \item Simbol $C$: $p(C) = \frac{1}{8} = 2^{-3} \implies \log_2(1/p(C)) = \log_2 8 = 3$ bit.
    \item Simbol $D$: $p(D) = \frac{1}{8} = 2^{-3} \implies \log_2(1/p(D)) = \log_2 8 = 3$ bit.
\end{itemize}
Maka nilai entropi adalah:
\begin{align*}
H(X) &= p(A)\log_2\left(\frac{1}{p(A)}\right) + p(B)\log_2\left(\frac{1}{p(B)}\right) + p(C)\log_2\left(\frac{1}{p(C)}\right) + p(D)\log_2\left(\frac{1}{p(D)}\right) \\
&= \left(\frac{1}{2} \times 1\right) + \left(\frac{1}{4} \times 2\right) + \left(\frac{1}{8} \times 3\right) + \left(\frac{1}{8} \times 3\right) \\
&= \frac{1}{2} + \frac{2}{4} + \frac{3}{8} + \frac{3}{8} = \frac{4}{8} + \frac{4}{8} + \frac{6}{8} = \frac{14}{8} = \frac{7}{4} = \mathbf{1.75\text{ bit}}.
\end{align*}

\noindent\textbf{Langkah 3: Perhitungan Bagian b (Entropi Maksimum $H_{\max}(X)$)}
Banyaknya simbol adalah $n = 4$. Berdasarkan Teorema Batas Entropi:
\[
H_{\max}(X) = \log_2(4) = \log_2(2^2) = \mathbf{2\text{ bit}}.
\]
Nilai maksimum 2 bit ini tercapai jika dan hanya jika seluruh simbol memiliki distribusi peluang seragam (\textit{uniform distribution}), yaitu:
\[
p(A) = p(B) = p(C) = p(D) = \frac{1}{4}.
\]

\noindent\textbf{Langkah 4: Perhitungan Bagian c (Efisiensi Informasi \& Redundansi)}
Efisiensi informasi $\eta$:
\[
\eta = \frac{H(X)}{H_{\max}(X)} = \frac{1.75}{2.00} = \mathbf{0.875} = \mathbf{87.5\%}.
\]
Redundansi sumber $R$:
\[
R = 1 - \eta = 1 - 0.875 = \mathbf{0.125} = \mathbf{12.5\%}.
\]

\vspace{0.4cm}

% =========================================================================
\subsection{Soal 2: Algoritma Huffman \& Teorema Batas Shannon}

\begin{soalmudah}{Soal 2 (Tingkat Mudah) -- Pengkodean Huffman Bebas-Prefiks \& Batas Shannon}
Diberikan alfabet sumber dengan 4 simbol $X = \{s_1, s_2, s_3, s_4\}$ yang memiliki distribusi peluang:
\[
p(s_1) = 0.40, \quad p(s_2) = 0.30, \quad p(s_3) = 0.20, \quad p(s_4) = 0.10.
\]
\textbf{Pertanyaan}:
\begin{enumerate}[label=\alph*)]
    \item Bangunlah pohon pengkodean Huffman (\textit{Huffman coding tree}) dan tentukan kata sandi biner (\textit{codeword}) untuk setiap simbol. Tunjukkan bahwa kode tersebut bersifat bebas-prefiks (\textit{prefix-free}).
    \item Hitunglah panjang rata-rata kode biner $l(f) = \sum p(s_i)|f(s_i)|$.
    \item Hitunglah nilai entropi $H(X)$ dan verifikasikan berlakunya Teorema Batas Pengkodean Sumber Shannon:
    \[
    H(X) \le l(f) < H(X) + 1.
    \]
\end{enumerate}
\end{soalmudah}

\noindent\textbf{Pembahasan Lengkap}:

\noindent\textbf{Langkah 1: Konstruksi Pohon Huffman (Bagian a)}
Urutan awal berdasarkan peluang menurun: $s_1 (0.40), s_2 (0.30), s_3 (0.20), s_4 (0.10)$.
\begin{itemize}
    \item \textit{Iterasi 1}: Gabungkan dua probabilitas terkecil, yaitu $s_3 (0.20)$ dan $s_4 (0.10)$ menjadi node gabungan $\{s_3, s_4\}$ dengan bobot $0.20 + 0.10 = 0.30$. Beri label bit: cabang $s_3 \to 1$, cabang $s_4 \to 0$.
    \item \textit{Iterasi 2}: Himpunan node saat ini: $s_1 (0.40), s_2 (0.30), \{s_3, s_4\} (0.30)$. Gabungkan dua terkecil: $s_2 (0.30)$ dan $\{s_3, s_4\} (0.30)$ menjadi $\{s_2, s_3, s_4\}$ dengan bobot $0.30 + 0.30 = 0.60$. Beri label bit: cabang $s_2 \to 0$, cabang $\{s_3, s_4\} \to 1$.
    \item \textit{Iterasi 3}: Gabungkan $s_1 (0.40)$ dan $\{s_2, s_3, s_4\} (0.60)$ menjadi akar (\textit{Root}) berbobot $1.00$. Beri label bit: cabang $s_1 \to 0$, cabang $\{s_2, s_3, s_4\} \to 1$.
\end{itemize}

\noindent\textbf{Tabel Hasil Pengkodean Huffman}:
\begin{center}
\begin{tabular}{ccccc}
\toprule
\textbf{Simbol} & \textbf{Peluang $p(s_i)$} & \textbf{Kata Sandi $f(s_i)$} & \textbf{Panjang $|f(s_i)|$} & $p(s_i) \cdot |f(s_i)|$ \\
\midrule
$s_1$ & $0.40$ & \texttt{0} & 1 bit & $0.40 \times 1 = 0.40$ \\
$s_2$ & $0.30$ & \texttt{10} & 2 bit & $0.30 \times 2 = 0.60$ \\
$s_3$ & $0.20$ & \texttt{111} & 3 bit & $0.20 \times 3 = 0.60$ \\
$s_4$ & $0.10$ & \texttt{110} & 3 bit & $0.10 \times 3 = 0.30$ \\
\bottomrule
\end{tabular}
\end{center}
\textit{Pemeriksaan Sifat Bebas-Prefiks}: Tidak ada kata sandi yang menjadi awalan dari kata sandi lain (\texttt{0} bukan awalan \texttt{10}, \texttt{111}, \texttt{110}; \texttt{10} bukan awalan \texttt{111} atau \texttt{110}; serta \texttt{111} bukan awalan \texttt{110}). Jadi kode bersifat \textbf{prefix-free}.

\noindent\textbf{Langkah 2: Panjang Rata-Rata Kode $l(f)$ (Bagian b)}
\[
l(f) = \sum_{i=1}^4 p(s_i)|f(s_i)| = 0.40 + 0.60 + 0.60 + 0.30 = \mathbf{1.90\text{ bit/simbol}}.
\]

\noindent\textbf{Langkah 3: Perhitungan Entropi $H(X)$ \& Verifikasi Teorema Batas (Bagian c)}
Gunakan nilai logaritma basis 2: $\log_2(0.4) \approx -1.3219$, $\log_2(0.3) \approx -1.7370$, $\log_2(0.2) \approx -2.3219$, $\log_2(0.1) \approx -3.3219$.
\begin{align*}
H(X) &= - \left[ 0.40(-1.3219) + 0.30(-1.7370) + 0.20(-2.3219) + 0.10(-3.3219) \right] \\
&= - \left[ -0.5288 - 0.5211 - 0.4644 - 0.3322 \right] = \mathbf{1.8465\text{ bit}}.
\end{align*}
Substitusikan ke Teorema Batas Shannon:
\[
1.8465 \le 1.90 < 1.8465 + 1 = 2.8465 \quad (\textbf{Terbukti Terpenuhi!})
\]

\vspace{0.4cm}

% =========================================================================
\subsection{Soal 3: Jarak Unicity Sandi Sederhana (Sandi Geser \& Sandi Affine)}

\begin{soalmudah}{Soal 3 (Tingkat Mudah) -- Jarak Unicity Sandi Geser \& Sandi Affine}
Misalkan teks bermakna ditulis dalam bahasa alami yang menggunakan 26 alfabet standar ($|\mathcal{P}| = 26$) dengan nilai entropi bahasa per karakter $H_L = 1.30\text{ bit/huruf}$.
\textbf{Pertanyaan}:
\begin{enumerate}[label=\alph*)]
    \item Hitunglah nilai redundansi bahasa alami tersebut ($R_L$).
    \item Hitunglah Jarak Unicity ($n_0$) untuk \textbf{Sandi Geser (\textit{Shift Cipher})}.
    \item Hitunglah Jarak Unicity ($n_0$) untuk \textbf{Sandi Affine} atas $\mathbb{Z}_{26}$.
    \item Jelaskan interpretasi operasional dari nilai $n_0$ tersebut bagi seorang kriptanalis yang melakukan serangan \textit{ciphertext-only}.
\end{enumerate}
\end{soalmudah}

\noindent\textbf{Pembahasan Lengkap}:

\noindent\textbf{Langkah 1: Perhitungan Redundansi Bahasa $R_L$ (Bagian a)}
Ukuran ruang alfabet: $|\mathcal{P}| = 26 \implies \log_2 26 = \frac{\ln 26}{\ln 2} \approx 4.7004\text{ bit}$.
Rumus redundansi bahasa:
\begin{equation}
R_L = 1 - \frac{H_L}{\log_2 |\mathcal{P}|} = 1 - \frac{1.30}{4.7004} \approx 1 - 0.27657 = \mathbf{0.72343} \approx \mathbf{72.34\%}.
\end{equation}

\noindent\textbf{Langkah 2: Jarak Unicity Sandi Geser (Bagian b)}
Pada Sandi Geser, kunci $K \in \mathbb{Z}_{26}$, sehingga ukuran ruang kunci $|\mathcal{K}| = 26$.
Maka $\log_2 |\mathcal{K}| = \log_2 26 \approx 4.7004\text{ bit}$.
Gunakan rumus Jarak Unicity:
\begin{equation}
\label{eq:unicity_formula}
n_0 \approx \frac{\log_2 |\mathcal{K}|}{R_L \log_2 |\mathcal{P}|} = \frac{\log_2 26}{R_L \log_2 26} = \frac{1}{R_L} = \frac{1}{0.72343} \approx 1.382 \approx \mathbf{2\text{ karakter}}.
\end{equation}

\noindent\textbf{Langkah 3: Jarak Unicity Sandi Affine (Bagian c)}
Pada Sandi Affine $e_{(a,b)}(x) = ax + b \bmod 26$, syarat keterbalikan menuntut $\gcd(a, 26) = 1$.
Banyaknya nilai $a$ yang relatif prima dengan 26 adalah $\phi(26) = \phi(2)\phi(13) = (1)(12) = 12$.
Banyaknya nilai pergeseran $b \in \mathbb{Z}_{26}$ adalah 26.
Total ruang kunci:
\[
|\mathcal{K}| = 12 \times 26 = 312 \implies \log_2 |\mathcal{K}| = \log_2(312) = \frac{\ln 312}{\ln 2} \approx 8.2854\text{ bit}.
\]
Substitusikan ke rumus:
\[
n_0 \approx \frac{8.2854}{0.72343 \times 4.7004} = \frac{8.2854}{3.4004} \approx 2.437 \approx \mathbf{3\text{ karakter}}.
\]

\noindent\textbf{Langkah 4: Interpretasi Kriptanalis (Bagian d)}
Nilai $n_0$ menyatakan panjang minimum teks sandi yang harus dicegat agar rata-rata jumlah kunci palsu berkurang menjadi 0.
\begin{itemize}
    \item Untuk Sandi Geser: Ciphertext sepanjang $\ge 2$ karakter sudah cukup bagi kriptanalis untuk menentukan kunci asli secara unik.
    \item Untuk Sandi Affine: Ciphertext sepanjang $\ge 3$ karakter secara teoretis menjamin hanya ada 1 kunci tunggal yang menghasilkan teks asal bermakna.
\end{itemize}

\vspace{0.4cm}

% =========================================================================
\subsection{Soal 4: Konsep Idempotensi \& Komposisi Cipher Sederhana}

\begin{soalmudah}{Soal 4 (Tingkat Mudah) -- Cipher Endomorfik \& Pembuktian Idempotensi Sandi Geser}
\textbf{Pertanyaan}:
\begin{enumerate}[label=\alph*)]
    \item Definisikan secara formal apa yang dimaksud dengan \textbf{sistem kripto endomorfik} dan \textbf{cipher idempoten}.
    \item Buktikan secara matematis bahwa Sandi Geser $S$ atas $\mathbb{Z}_{26}$ bersifat idempoten, yaitu $S^2 = S \times S = S$.
    \item Mengapa enkripsi ganda (\textit{double encryption}) menggunakan Sandi Geser dengan dua kunci berbeda $K_1$ dan $K_2$ sama sekali tidak meningkatkan kekuatan keamanan terhadap serangan \textit{brute-force}?
\end{enumerate}
\end{soalmudah}

\noindent\textbf{Pembahasan Lengkap}:

\noindent\textbf{Langkah 1: Definisi Formal (Bagian a)}
\begin{itemize}
    \item \textbf{Cipher Endomorfik}: Suatu sistem kripto $(\mathcal{P}, \mathcal{C}, \mathcal{K}, \mathcal{E}, \mathcal{D})$ disebut endomorfik jika ruang teks sandi sama dengan ruang teks asal, yaitu $\mathcal{C} = \mathcal{P}$.
    \item \textbf{Cipher Idempoten}: Suatu sistem kripto endomorfik $S$ dikatakan idempoten jika perkalian komposisi sistem dengan dirinya sendiri menghasilkan sistem yang identik:
    \begin{equation}
    S^2 = S \times S = S.
    \end{equation}
\end{itemize}

\noindent\textbf{Langkah 2: Pembuktian Matematis Idempotensi Sandi Geser (Bagian b)}
Misalkan $S = (\mathbb{Z}_{26}, \mathbb{Z}_{26}, \mathbb{Z}_{26}, \mathcal{E}, \mathcal{D})$ dengan fungsi enkripsi $e_K(x) = (x + K) \bmod 26$ dan peluang kunci seragam $p(K) = \frac{1}{26}$.
Tinjau perkalian sistem $S_1 \times S_2$:
\begin{enumerate}
    \item Ruang kunci gabungan adalah pasangan $(K_1, K_2) \in \mathbb{Z}_{26} \times \mathbb{Z}_{26}$ dengan peluang $p(K_1, K_2) = \frac{1}{26} \times \frac{1}{26} = \frac{1}{676}$.
    \item Fungsi enkripsi gabungan adalah komposisi berantai:
    \[
    e_{(K_1, K_2)}(x) = e_{K_2}(e_{K_1}(x)) = ((x + K_1) + K_2) \bmod 26 = (x + (K_1 + K_2 \bmod 26)) \bmod 26.
    \]
    \item Definisikan $K_3 = (K_1 + K_2) \bmod 26$. Karena $(\mathbb{Z}_{26}, +)$ adalah grup hingga, maka $K_3 \in \mathbb{Z}_{26}$.
    \item Untuk setiap nilai $k \in \mathbb{Z}_{26}$, terdapat tepat 26 pasangan $(K_1, K_2)$ yang memenuhi $(K_1 + K_2) \equiv k \pmod{26}$, yaitu pasangan $(K_1, (k - K_1) \bmod 26)$ untuk $K_1 = 0, 1, \dots, 25$.
    \item Probabilitas munculnya kunci efektif $K_3 = k$ adalah:
    \[
    P(K_3 = k) = \sum_{K_1=0}^{25} p(K_1, (k - K_1) \bmod 26) = 26 \times \frac{1}{676} = \frac{1}{26}.
    \]
\end{enumerate}
Karena operasi enkripsi gabungan $e_{K_3}(x) = (x + K_3) \bmod 26$ dan distribusi peluang kuncinya $P(K_3 = k) = \frac{1}{26}$ identik dengan Sandi Geser semula $S$, maka terbukti secara formal:
\[
\mathbf{S \times S = S \quad (\text{Idempoten})}.
\]

\noindent\textbf{Langkah 3: Analisis Keamanan Enkripsi Ganda (Bagian c)}
Karena $S^2 = S$, enkripsi ganda dengan dua kunci geser $K_1, K_2$ menghasilkan transformasi yang ekuivalen dengan enkripsi tunggal berkunci $K_3 = (K_1 + K_2) \bmod 26$. Penyerang yang melakukan pencarian kunci secara menyeluruh (\textit{exhaustive key search}) tidak perlu mencoba $26^2 = 676$ pasangan, melainkan hanya perlu menguji 26 kemungkinan nilai $K_3$. Kompleksitas serangan tetap $O(26)$, sehingga \textbf{tidak ada peningkatan keamanan sama sekali}.

\vspace{0.6cm}

% =========================================================================
\section{Tingkat 2: Tingkat Sedang (Soal 5 -- 8)}
% =========================================================================

\subsection{Soal 5: Uji Keamanan Sempurna Matriks Kripto \texorpdfstring{$3 \times 3$}{3x3}}

\begin{soalsedang}{Soal 5 (Tingkat Sedang) -- Model Probabilitas \& Uji Keamanan Sempurna Matriks $3 \times 3$}
Diberikan sebuah sistem kriptografi $(\mathcal{P}, \mathcal{C}, \mathcal{K}, \mathcal{E}, \mathcal{D})$ dengan spesifikasi:
\begin{itemize}
    \item Ruang plaintext: $\mathcal{P} = \{m_1, m_2, m_3\}$ dengan distribusi $p_{\mathcal{P}}(m_1) = \frac{1}{2}, \; p_{\mathcal{P}}(m_2) = \frac{1}{4}, \; p_{\mathcal{P}}(m_3) = \frac{1}{4}$.
    \item Ruang kunci: $\mathcal{K} = \{k_1, k_2, k_3\}$ dengan probabilitas seragam $p_{\mathcal{K}}(k_1) = p_{\mathcal{K}}(k_2) = p_{\mathcal{K}}(k_3) = \frac{1}{3}$.
    \item Ruang ciphertext: $\mathcal{C} = \{c_1, c_2, c_3\}$.
    \item Matriks Enkripsi $e_k(m)$:
    \begin{center}
    \begin{tabular}{c|ccc}
    \toprule
    Kunci & Plaintext $m_1$ & Plaintext $m_2$ & Plaintext $m_3$ \\
    \midrule
    $k_1$ & $c_1$ & $c_2$ & $c_3$ \\
    $k_2$ & $c_2$ & $c_3$ & $c_1$ \\
    $k_3$ & $c_3$ & $c_1$ & $c_2$ \\
    \bottomrule
    \end{tabular}
    \end{center}
\end{itemize}
\textbf{Pertanyaan}:
\begin{enumerate}[label=\alph*)]
    \item Hitunglah distribusi peluang marginal ciphertext $p_{\mathcal{C}}(c_1), p_{\mathcal{C}}(c_2), p_{\mathcal{C}}(c_3)$.
    \item Hitunglah peluang a posteriori $p_{\mathcal{P}}(m_i | c_j)$ untuk setiap pasangan $(m_i, c_j)$ dengan aturan Bayes.
    \item Tentukan apakah sistem ini memiliki \textbf{keamanan sempurna (\textit{perfect secrecy})} dan buktikan keterkaitannya dengan Teorema Karakterisasi Shannon 3.
\end{enumerate}
\end{soalsedang}

\noindent\textbf{Pembahasan Lengkap}:

\noindent\textbf{Langkah 1: Menghitung Peluang Marginal $p_{\mathcal{C}}(c_j)$ (Bagian a)}
Gunakan rumus probabilitas total $p_{\mathcal{C}}(y) = \sum_{k \in \mathcal{K}} p_{\mathcal{K}}(k) \cdot p_{\mathcal{P}}(d_k(y))$:
\begin{itemize}
    \item Ciphertext $c_1$ dihasilkan oleh: $(k_1, m_1)$, $(k_2, m_3)$, $(k_3, m_2)$.
    \[
    p_{\mathcal{C}}(c_1) = \frac{1}{3} p(m_1) + \frac{1}{3} p(m_3) + \frac{1}{3} p(m_2) = \frac{1}{3}\left(\frac{1}{2} + \frac{1}{4} + \frac{1}{4}\right) = \frac{1}{3}(1) = \mathbf{\frac{1}{3}}.
    \]
    \item Ciphertext $c_2$ dihasilkan oleh: $(k_1, m_2)$, $(k_2, m_1)$, $(k_3, m_3)$.
    \[
    p_{\mathcal{C}}(c_2) = \frac{1}{3} p(m_2) + \frac{1}{3} p(m_1) + \frac{1}{3} p(m_3) = \frac{1}{3}\left(\frac{1}{4} + \frac{1}{2} + \frac{1}{4}\right) = \frac{1}{3}(1) = \mathbf{\frac{1}{3}}.
    \]
    \item Ciphertext $c_3$ dihasilkan oleh: $(k_1, m_3)$, $(k_2, m_2)$, $(k_3, m_1)$.
    \[
    p_{\mathcal{C}}(c_3) = \frac{1}{3} p(m_3) + \frac{1}{3} p(m_2) + \frac{1}{3} p(m_1) = \frac{1}{3}\left(\frac{1}{4} + \frac{1}{4} + \frac{1}{2}\right) = \frac{1}{3}(1) = \mathbf{\frac{1}{3}}.
    \]
\end{itemize}
Distribusi marginal ciphertext berdistribusi seragam: $p_{\mathcal{C}}(c_j) = \frac{1}{3}$ untuk seluruh $j \in \{1, 2, 3\}$.

\noindent\textbf{Langkah 2: Menghitung Peluang Bersyarat $p(m_i | c_j)$ (Bagian b)}
Gunakan Aturan Bayes $p(m_i | c_j) = \frac{p(m_i) p(c_j | m_i)}{p(c_j)}$.
Perhatikan bahwa untuk setiap pesan $m_i$ dan sembarang ciphertext $c_j$, terdapat tepat 1 kunci unik yang memetakan $m_i \to c_j$, sehingga $p(c_j | m_i) = p(k) = \frac{1}{3}$.
Maka:
\[
p(m_i | c_j) = \frac{p(m_i) \cdot \frac{1}{3}}{\frac{1}{3}} = p(m_i).
\]
Secara eksplisit:
\begin{itemize}
    \item Untuk $m_1$: $p(m_1 | c_1) = p(m_1 | c_2) = p(m_1 | c_3) = \mathbf{\frac{1}{2}} = p(m_1)$.
    \item Untuk $m_2$: $p(m_2 | c_1) = p(m_2 | c_2) = p(m_2 | c_3) = \mathbf{\frac{1}{4}} = p(m_2)$.
    \item Untuk $m_3$: $p(m_3 | c_1) = p(m_3 | c_2) = p(m_3 | c_3) = \mathbf{\frac{1}{4}} = p(m_3)$.
\end{itemize}

\noindent\textbf{Langkah 3: Kesimpulan Keamanan Sempurna \& Teorema Shannon (Bagian c)}
Karena $p(m_i | c_j) = p(m_i)$ berlaku untuk \textbf{seluruh} $m_i \in \mathcal{P}$ dan seluruh $c_j \in \mathcal{C}$, maka sistem ini \textbf{MEMILIKI KEAMANAN SEMPURNA (\textit{Perfect Secrecy})}.
\textit{Kesesuaian dengan Teorema Shannon 3}:
\begin{enumerate}
    \item Ukuran ketiga ruang sama: $|\mathcal{K}| = |\mathcal{C}| = |\mathcal{P}| = 3$.
    \item Kunci dipilih seragam: $p_{\mathcal{K}}(k) = \frac{1}{3} = \frac{1}{|\mathcal{K}|}$.
    \item Setiap baris/kolom matriks enkripsi merupakan permutasi lengkap (untuk setiap pasangan $(m_i, c_j)$ ada tepat satu kunci $k$).
\end{enumerate}

\vspace{0.4cm}

% =========================================================================
\subsection{Soal 6: Entropi Bersama, Bersyarat, \& Key Equivocation}

\begin{soalsedang}{Soal 6 (Tingkat Sedang) -- Perhitungan Key Equivocation $H(K|C)$ \& Entropi Bersyarat}
Menggunakan sistem kriptografi yang didefinisikan pada Soal 5 di atas:
\textbf{Pertanyaan}:
\begin{enumerate}[label=\alph*)]
    \item Hitunglah nilai entropi plaintext $H(P)$, entropi kunci $H(K)$, dan entropi ciphertext $H(C)$.
    \item Hitunglah \textit{Key Equivocation} $H(K|C)$ dengan menggunakan Rumus Teorema Shannon:
    \[
    H(K|C) = H(K) + H(P) - H(C).
    \]
    \item Hitunglah entropi bersyarat teks asal $H(P|C)$ dan jelaskan mengapa nilainya sama persis dengan $H(K|C)$ pada sistem ini.
    \item Jelaskan makna teoretis dari hasil $H(P|C) = H(P)$ kaitannya dengan sifat \textit{perfect secrecy}.
\end{enumerate}
\end{soalsedang}

\noindent\textbf{Pembahasan Lengkap}:

\noindent\textbf{Langkah 1: Perhitungan $H(P), H(K), H(C)$ (Bagian a)}
\begin{enumerate}
    \item \textbf{Entropi Plaintext $H(P)$}:
    \begin{align*}
    H(P) &= - \left[ \frac{1}{2}\log_2\left(\frac{1}{2}\right) + \frac{1}{4}\log_2\left(\frac{1}{4}\right) + \frac{1}{4}\log_2\left(\frac{1}{4}\right) \right] \\
    &= - \left[ \frac{1}{2}(-1) + \frac{1}{4}(-2) + \frac{1}{4}(-2) \right] = \frac{1}{2} + \frac{1}{2} + \frac{1}{2} = \mathbf{1.5\text{ bit}}.
    \end{align*}
    \item \textbf{Entropi Kunci $H(K)$}:
    Karena 3 kunci berdistribusi seragam:
    \[
    H(K) = \log_2 3 \approx \mathbf{1.58496\text{ bit}}.
    \]
    \item \textbf{Entropi Ciphertext $H(C)$}:
    Dari Soal 5, distribusi marginal ciphertext seragam ($p(c_j) = 1/3$), sehingga:
    \[
    H(C) = \log_2 3 \approx \mathbf{1.58496\text{ bit}}.
    \]
\end{enumerate}

\noindent\textbf{Langkah 2: Perhitungan Key Equivocation $H(K|C)$ (Bagian b)}
Gunakan formula Shannon:
\begin{align*}
H(K|C) &= H(K) + H(P) - H(C) \\
&= \log_2 3 + 1.5 - \log_2 3 = \mathbf{1.5\text{ bit}}.
\end{align*}

\noindent\textbf{Langkah 3: Perhitungan $H(P|C)$ \& Kesamaan Nilai (Bagian c)}
Definisi entropi bersyarat:
\[
H(P|C) = \sum_{j=1}^3 p(c_j) H(P|c_j).
\]
Karena $p(m_i|c_j) = p(m_i)$ untuk setiap $j$, maka distribusi bersyarat $P|c_j$ identik dengan distribusi apriori $P$.
Maka $H(P|c_j) = H(P) = 1.5\text{ bit}$ untuk setiap $j$, sehingga:
\[
H(P|C) = \frac{1}{3}(1.5) + \frac{1}{3}(1.5) + \frac{1}{3}(1.5) = \mathbf{1.5\text{ bit}}.
\]
\textit{Alasan Kesamaan $H(K|C) = H(P|C)$}: Karena untuk ciphertext $c_j$ yang diamati, setiap pilihan kunci $k$ berkorespodensi satu-ke-satu dengan tepat satu plaintext $m = d_k(c_j)$. Mengetahui kunci setara dengan mengetahui pesan asal, sehingga ketidakpastian kunci $H(K|C)$ sama dengan ketidakpastian pesan $H(P|C)$.

\noindent\textbf{Langkah 4: Makna Teoretis Kaitannya dengan Keamanan Sempurna (Bagian d)}
Kondisi $H(P|C) = H(P)$ menunjukkan bahwa pengamatan terhadap ciphertext $C$ sama sekali tidak mengurangi ketidakpastian penyerang mengenai plaintext $P$. Informasi saling (\textit{Mutual Information}) bernilai nol:
\[
I(P; C) = H(P) - H(P|C) = 1.5 - 1.5 = 0\text{ bit}.
\]
Ini adalah karakteristik informasi murni dari sistem dengan \textbf{keamanan sempurna}.

\vspace{0.4cm}

% =========================================================================
\subsection{Soal 7: Jarak Unicity Sandi Vigenère \& Sandi Permutasi}

\begin{soalsedang}{Soal 7 (Tingkat Sedang) -- Analisis Jarak Unicity Sandi Polialfabetik vs Sandi Permutasi}
Teks bahasa Inggris memiliki redundansi $R_L = 0.75$ pada alfabet berukuran $|\mathcal{P}| = 26$ ($\log_2 26 \approx 4.7004\text{ bit}$).
\textbf{Pertanyaan}:
\begin{enumerate}[label=\alph*)]
    \item Turunkan rumus umum Jarak Unicity $n_0$ untuk \textbf{Sandi Vigen\`ere} dengan panjang kunci periode $m$. Hitung nilai numerik $n_0$ untuk $m = 8$.
    \item Sebuah \textbf{Sandi Permutasi} membagi pesan menjadi blok-blok berukuran $m = 8$ huruf dan mempermutasikan urutan huruf di dalam setiap blok menggunakan suatu permutasi $\pi \in S_8$. Tentukan ukuran ruang kunci $|\mathcal{K}|$ dan hitung nilai $n_0$ untuk sandi permutasi ini.
    \item Bandingkan kedua nilai $n_0$ tersebut dan jelaskan mengapa Sandi Permutasi memiliki jarak unicity yang jauh lebih kecil daripada Sandi Vigen\`ere meskipun keduanya bekerja pada panjang blok yang sama ($m = 8$).
\end{enumerate}
\end{soalsedang}

\noindent\textbf{Pembahasan Lengkap}:

\noindent\textbf{Langkah 1: Rumus Umum \& Perhitungan Sandi Vigenère (Bagian a)}
Ruang kunci Sandi Vigenère periode $m$ terdiri dari seluruh untai $m$-huruf alfabet: $|\mathcal{K}| = 26^m$.
Maka pembilang entropi kuncinya adalah:
\[
\log_2 |\mathcal{K}| = \log_2(26^m) = m \log_2 26.
\]
Substitusikan ke rumus jarak unicity:
\begin{equation}
n_0 \approx \frac{m \log_2 26}{R_L \log_2 26} = \mathbf{\frac{m}{R_L}}.
\end{equation}
Untuk $R_L = 0.75 = \frac{3}{4}$ dan periode $m = 8$:
\[
n_0 \approx \frac{8}{0.75} = \frac{8 \times 4}{3} = \frac{32}{3} \approx 10.67 \approx \mathbf{11\text{ karakter}}.
\]

\noindent\textbf{Langkah 2: Perhitungan Sandi Permutasi Blok $m = 8$ (Bagian b)}
Kunci sandi permutasi adalah elemen grup simetri $S_8$, sehingga:
\[
|\mathcal{K}| = 8! = 8 \times 7 \times 6 \times 5 \times 4 \times 3 \times 2 \times 1 = 40.320.
\]
Hitung entropi kunci:
\[
\log_2 |\mathcal{K}| = \log_2(40.320) = \frac{\ln 40320}{\ln 2} \approx 15.2992\text{ bit}.
\]
Hitung jarak unicity:
\[
n_0 \approx \frac{15.2992}{R_L \log_2 |\mathcal{P}|} = \frac{15.2992}{0.75 \times 4.7004} = \frac{15.2992}{3.5253} \approx 4.3398 \approx \mathbf{5\text{ karakter}}.
\]

\noindent\textbf{Langkah 3: Perbandingan \& Analisis Kriptografis (Bagian c)}
\begin{itemize}
    \item Ukuran ruang kunci Vigenère ($m=8$): $|\mathcal{K}| = 26^8 \approx 2.088 \times 10^{11}$ ($\log_2 |\mathcal{K}| \approx 37.60\text{ bit}$).
    \item Ukuran ruang kunci Permutasi ($m=8$): $|\mathcal{K}| = 8! = 40.320$ ($\log_2 |\mathcal{K}| \approx 15.30\text{ bit}$).
\end{itemize}
\textbf{Alasan}: Sandi permutasi hanya mengacak posisi karakter tanpa mengubah identitas huruf, sehingga ruang kuncinya ($8! = 40.320$) jauh lebih sempit dibandingkan ruang kunci substitusi polialfabetik ($26^8$). Akibatnya, informasi redundansi bahasa alami mengeliminasi kunci palsu jauh lebih cepat pada sandi permutasi ($n_0 \approx 5$ huruf, yaitu kurang dari 1 blok!) dibandingkan Vigenère ($n_0 \approx 11$ huruf).

\vspace{0.4cm}

% =========================================================================
\subsection{Soal 8: Komutativitas \& Idempotensi Komposisi Cipher Endomorfik}

\begin{soalsedang}{Soal 8 (Tingkat Sedang) -- Pembuktian $M \times S = S \times M = \text{Affine}$ \& Non-Komutatif}
Diberikan sistem Sandi Multiplikatif $M$ ($e_a(x) = ax \bmod 26$, dengan $\gcd(a, 26) = 1$) dan Sandi Geser $S$ ($e_b(x) = x + b \bmod 26$, dengan $b \in \mathbb{Z}_{26}$).
\textbf{Pertanyaan}:
\begin{enumerate}[label=\alph*)]
    \item Buktikan secara aljabar bahwa $M \times S = S \times M = \text{Sandi Affine}$.
    \item Tunjukkan bahwa Sandi Affine bersifat idempoten ($\text{Affine} \times \text{Affine} = \text{Affine}$).
    \item Misalkan $S_{\text{sub}}$ adalah sandi substitusi monoalfabetik dan $S_{\text{trans}}$ adalah sandi permutasi transposisi kolom. Jelaskan mengapa komposisi keduanya $S_{\text{sub}} \times S_{\text{trans}}$ secara umum \textbf{tidak bersifat komutatif} dan \textbf{tidak bersifat idempoten}.
\end{enumerate}
\end{soalsedang}

\noindent\textbf{Pembahasan Lengkap}:

\noindent\textbf{Langkah 1: Pembuktian Aljabar Komutativitas $M \times S = S \times M$ (Bagian a)}
\begin{itemize}
    \item \textbf{Komposisi $M \times S$}: Terapkan $M$ terlebih dahulu dengan kunci $a \in \mathbb{Z}_{26}^*$, lalu $S$ dengan kunci $b \in \mathbb{Z}_{26}$:
    \[
    e_{(a, b)}(x) = e_b(e_a(x)) = (ax + b) \bmod 26.
    \]
    Ini persis merupakan rumus enkripsi Sandi Affine berkunci $(a, b) \in \mathbb{Z}_{26}^* \times \mathbb{Z}_{26}$ berpeluang $p = \frac{1}{12 \times 26} = \frac{1}{312}$.
    \item \textbf{Komposisi $S \times M$}: Terapkan $S$ terlebih dahulu dengan kunci $b \in \mathbb{Z}_{26}$, lalu $M$ dengan kunci $a \in \mathbb{Z}_{26}^*$:
    \[
    e_{(b, a)}(x) = e_a(e_b(x)) = a(x + b) \bmod 26 = (ax + ab) \bmod 26.
    \]
    Definisikan $b' = ab \bmod 26$. Karena $\gcd(a, 26) = 1$, perkalian dengan $a$ modulo 26 merupakan pemetaan bijektif (permutasi) pada $\mathbb{Z}_{26}$.
    Jika $b$ berdistribusi seragam pada $\mathbb{Z}_{26}$ dengan peluang $\frac{1}{26}$, maka $b'$ juga berdistribusi seragam pada $\mathbb{Z}_{26}$ dengan peluang $\frac{1}{26}$.
\end{itemize}
Karena kedua komposisi menghasilkan himpunan operasi enkripsi yang sama ($x \mapsto ax + c \bmod 26$) dengan distribusi probabilitas kunci seragam yang identik ($\frac{1}{312}$), maka terbukti:
\[
\mathbf{M \times S = S \times M = \text{Sandi Affine}}.
\]

\noindent\textbf{Langkah 2: Pembuktian Idempotensi Sandi Affine (Bagian b)}
Tinjau komposisi dua Sandi Affine $A_1$ berkunci $(a_1, b_1)$ dan $A_2$ berkunci $(a_2, b_2)$:
\begin{align*}
e_{(a_2, b_2)}(e_{(a_1, b_1)}(x)) &= a_2(a_1 x + b_1) + b_2 \bmod 26 \\
&= (a_2 a_1)x + (a_2 b_1 + b_2) \bmod 26.
\end{align*}
Definisikan $a_3 = (a_2 a_1) \bmod 26$ dan $b_3 = (a_2 b_1 + b_2) \bmod 26$.
Karena $\gcd(a_1, 26) = 1$ dan $\gcd(a_2, 26) = 1$, maka $\gcd(a_3, 26) = 1$. Jelas $b_3 \in \mathbb{Z}_{26}$.
Pasangan $(a_3, b_3)$ merupakan kunci sah Sandi Affine. Maka $\text{Affine} \times \text{Affine} = \text{Affine}$ (Idempoten).

\noindent\textbf{Langkah 3: Analisis $S_{\text{sub}} \times S_{\text{trans}}$ (Bagian c)}
\begin{enumerate}
    \item \textbf{Non-Komutatif}: $S_{\text{sub}}$ mengganti nilai simbol alfabet pada posisi tertentu ($x_i \mapsto \pi(x_i)$), sedangkan $S_{\text{trans}}$ menukar posisi spasial huruf dalam blok ($i \mapsto \sigma(i)$). Jika dilakukan transposisi sebelum substitusi polialfabetik/berblok, karakter yang bertemu dengan kunci tertentu berubah total dibandingkan jika substitusi dilakukan lebih dulu.
    \item \textbf{Non-Idempoten}: $(S_{\text{sub}} \times S_{\text{trans}})^2 \neq S_{\text{sub}} \times S_{\text{trans}}$ karena operasi substitusi non-linear yang diselingi permutasi posisi menghasilkan fungsi pemetaan global dengan derajat kompleksitas yang berlipat ganda, yang tidak dapat direduksi menjadi satu lapis substitusi-transposisi tunggal. Prinsip inilah yang menjadi dasar iterasi ronde SPN modern.
\end{enumerate}

\vspace{0.6cm}

% =========================================================================
\section{Tingkat 3: Tingkat Sulit (Soal 9 -- 12)}
% =========================================================================

\subsection{Soal 9: Pembuktian Formal Syarat Shannon \texorpdfstring{$|\mathcal{K}| \ge |\mathcal{C}| \ge |\mathcal{P}|$}{|K| >= |C| >= |P|}}

\begin{soalsulit}{Soal 9 (Tingkat Sulit) -- Teorema Ketaksamaan Ukuran Ruang Kunci Shannon}
Misalkan $(\mathcal{P}, \mathcal{C}, \mathcal{K}, \mathcal{E}, \mathcal{D})$ adalah sistem kriptografi yang memiliki \textbf{keamanan sempurna (\textit{perfect secrecy})}, dengan asumsi seluruh $x \in \mathcal{P}$ memiliki peluang $p_{\mathcal{P}}(x) > 0$ dan seluruh $y \in \mathcal{C}$ memiliki peluang $p_{\mathcal{C}}(y) > 0$.
\textbf{Pertanyaan}:
\begin{enumerate}[label=\alph*)]
    \item Buktikan secara formal bahwa untuk setiap pasangan $(x, y) \in \mathcal{P} \times \mathcal{C}$, terdapat sedikitnya satu kunci $K \in \mathcal{K}$ sedemikian sehingga $e_K(x) = y$.
    \item Berdasarkan hasil (a), buktikan bahwa harus berlaku ketaksamaan:
    \[
    |\mathcal{K}| \ge |\mathcal{C}|.
    \]
    \item Buktikan bahwa $|\mathcal{C}| \ge |\mathcal{P}|$ sehingga disimpulkan rantai ketaksamaan fundamental Shannon:
    \[
    |\mathcal{K}| \ge |\mathcal{C}| \ge |\mathcal{P}|.
    \]
    \item Jelaskan implikasi praktis teorema ini dalam desain kriptografi modern, khususnya mengapa \textit{One-Time Pad} (OTP) tidak dapat diterapkan untuk enkripsi data massal (misal hard drive atau streaming video gigabyte).
\end{enumerate}
\end{soalsulit}

\noindent\textbf{Pembahasan Lengkap}:

\noindent\textbf{Langkah 1: Pembuktian Eksistensi Kunci untuk Setiap Pasangan $(x, y)$ (Bagian a)}
\begin{proof}
Berdasarkan definisi keamanan sempurna (\textit{perfect secrecy}):
\begin{equation}
p_{\mathcal{P}}(x|y) = p_{\mathcal{P}}(x) \quad \text{untuk setiap } x \in \mathcal{P}, \, y \in \mathcal{C}.
\end{equation}
Diketahui bahwa $p_{\mathcal{P}}(x) > 0$. Akibatnya, $p_{\mathcal{P}}(x|y) > 0$.
Berdasarkan Teorema Bayes:
\begin{equation}
p_{\mathcal{P}}(x|y) = \frac{p_{\mathcal{P}}(x) \cdot p_{\mathcal{C}}(y|x)}{p_{\mathcal{C}}(y)} > 0 \implies p_{\mathcal{C}}(y|x) > 0.
\end{equation}
Di sisi lain, dari rumus peluang bersyarat sistem kripto:
\begin{equation}
p_{\mathcal{C}}(y|x) = \sum_{\{K \in \mathcal{K} \,:\, e_K(x) = y\}} p_{\mathcal{K}}(K).
\end{equation}
Agar jumlahan di atas bernilai $> 0$, himpunan indeks kunci $\{K \in \mathcal{K} : e_K(x) = y\}$ tidak boleh merupakan himpunan kosong.
Artinya, harus terdapat sekurang-kurangnya satu kunci $K \in \mathcal{K}$ sedemikian sehingga $e_K(x) = y$.
\end{proof}

\noindent\textbf{Langkah 2: Pembuktian $|\mathcal{K}| \ge |\mathcal{C}|$ (Bagian b)}
\begin{proof}
Pilih dan tetapkan satu elemen plaintext sebarang $x_0 \in \mathcal{P}$.
Berdasarkan hasil Bagian (a), untuk setiap $y \in \mathcal{C}$, terdapat sedikitnya satu kunci $K_y \in \mathcal{K}$ sedemikian sehingga:
\[
e_{K_y}(x_0) = y.
\]
Karena fungsi enkripsi $e_K$ adalah fungsi deterministik berharga tunggal, jika $y_1 \neq y_2$, maka tidak mungkin kunci yang sama memetakan $x_0$ ke dua bayangan yang berbeda:
\[
e_{K_{y_1}}(x_0) = y_1 \neq y_2 = e_{K_{y_2}}(x_0) \implies K_{y_1} \neq K_{y_2}.
\]
Dengan demikian, pemetaan $y \mapsto K_y$ mendefinisikan sebuah fungsi injektif dari himpunan $\mathcal{C}$ ke himpunan $\mathcal{K}$.
Berdasarkan prinsip dasar teori himpunan, keberadaan fungsi injektif $\mathcal{C} \hookrightarrow \mathcal{K}$ membuktikan bahwa:
\[
\mathbf{|\mathcal{K}| \ge |\mathcal{C}|}.
\]
\end{proof}

\noindent\textbf{Langkah 3: Pembuktian $|\mathcal{C}| \ge |\mathcal{P}|$ \& Rantai Ketaksamaan (Bagian c)}
\begin{proof}
Pilih satu kunci sebarang $K_0 \in \mathcal{K}$.
Karena proses dekripsi harus selalu menghasilkan teks asal semula tanpa ambigu, yaitu $d_{K_0}(e_{K_0}(x)) = x$ untuk setiap $x \in \mathcal{P}$, maka fungsi enkripsi $e_{K_0}: \mathcal{P} \to \mathcal{C}$ harus merupakan fungsi injektif.
Keberadaan fungsi injektif dari $\mathcal{P}$ ke $\mathcal{C}$ membuktikan bahwa:
\[
|\mathcal{C}| \ge |\mathcal{P}|.
\]
Menggabungkan ketaksamaan pada Bagian (b) dan Bagian (c), diperoleh rantai ketaksamaan fundamental Shannon:
\[
\mathbf{|\mathcal{K}| \ge |\mathcal{C}| \ge |\mathcal{P}|}.
\]
\end{proof}

\noindent\textbf{Langkah 4: Implikasi Praktis (Bagian d)}
Teorema Shannon membuktikan batasan teoritis yang mutlak: \textbf{untuk mencapai keamanan sempurna tanpa syarat, ukuran kunci minimal harus sama panjangnya dengan ukuran seluruh pesan yang ingin dikirimkan} ($H(K) \ge H(P)$).
\begin{itemize}
    \item Jika ingin mengenkripsi file sebesar 10 GB dengan OTP, pihak pengirim dan penerima harus terlebih dahulu menyepakati dan mendistribusikan kunci acak sejati sebesar 10 GB melalui saluran rahasia.
    \item Masalah distribusi kunci (\textit{key distribution problem}) sebesar ini sama sulitnya dengan mengirim pesan aslinya secara aman.
    \item Oleh karena itu, kriptografi praktis modern mengorbankan \textit{unconditional security} dan beralih ke \textbf{\textit{computational security}} (seperti AES dengan kunci pendek 128/256 bit atau RSA), di mana keamanan dijamin oleh ketidakmungkinan komputasi penyerang dalam waktu polinomial.
\end{itemize}

\vspace{0.4cm}

% =========================================================================
\subsection{Soal 10: Pembuktian Ketaksamaan Subaditivitas Entropi via Jensen}

\begin{soalsulit}{Soal 10 (Tingkat Sulit) -- Kecekungan $\log_2(t)$ \& Pembuktian Formal Subaditivitas $H(X,Y)$}
\textbf{Pertanyaan}:
\begin{enumerate}[label=\alph*)]
    \item Tuliskan pernyataan formal \textbf{Ketaksamaan Jensen} untuk fungsi cekung (\textit{concave function}) dan buktikan bahwa fungsi $f(t) = \log_2 t$ adalah fungsi cekung tegas (\textit{strictly concave}) pada interval $(0, \infty)$.
    \item Buktikan secara analitik bahwa untuk dua variabel acak diskret $X$ dan $Y$, berlaku sifat subaditivitas entropi:
    \[
    H(X, Y) \le H(X) + H(Y).
    \]
    \textit{Petunjuk}: Bentuklah ekspresi $H(X, Y) - H(X) - H(Y)$ sebagai jumlahan tunggal dengan logaritma rasio, lalu terapkan Ketaksamaan Jensen.
    \item Tentukan kondisi perlu dan cukup agar tanda kesamaan $H(X, Y) = H(X) + H(Y)$ tercapai.
    \item Buktikan bahwa observasi bersyarat tidak pernah menambah rata-rata entropi, yaitu $H(X|Y) \le H(X)$.
\end{enumerate}
\end{soalsulit}

\noindent\textbf{Pembahasan Lengkap}:

\noindent\textbf{Langkah 1: Ketaksamaan Jensen \& Kecekungan $f(t) = \log_2 t$ (Bagian a)}
\begin{itemize}
    \item \textbf{Teorema Ketaksamaan Jensen}: Misalkan $f: I \to \mathbb{R}$ adalah fungsi cekung pada interval terbuka $I$. Untuk sembarang bobot real $a_i > 0$ dengan $\sum_{i=1}^m a_i = 1$ dan untuk sembarang $t_i \in I$:
    \begin{equation}
    \label{eq:jensen_app}
    \sum_{i=1}^m a_i f(t_i) \le f\left( \sum_{i=1}^m a_i t_i \right).
    \end{equation}
    Jika $f$ cekung tegas, kesamaan berlaku jika dan hanya jika $t_1 = t_2 = \dots = t_m$.
    \item \textbf{Bukti Kecekungan Tegas}:
    Fungsi $f(t) = \log_2 t = \frac{\ln t}{\ln 2}$. Turunan pertama dan keduanya adalah:
    \[
    f'(t) = \frac{1}{t \ln 2}, \quad f''(t) = -\frac{1}{t^2 \ln 2}.
    \]
    Karena $\ln 2 > 0$ dan $t^2 > 0$ untuk setiap $t \in (0, \infty)$, maka $f''(t) < 0$ untuk seluruh $t \in (0, \infty)$. Karena turunan kedua senantiasa negatif, terbukti $f(t) = \log_2 t$ adalah \textbf{fungsi cekung tegas}.
\end{itemize}

\noindent\textbf{Langkah 2: Pembuktian Subaditivitas $H(X, Y) \le H(X) + H(Y)$ (Bagian b)}
\begin{proof}
Misalkan distribusi bersama $P(X = x_i, Y = y_j) = r_{ij} > 0$, distribusi marginal $P(X = x_i) = p_i = \sum_j r_{ij}$, dan $P(Y = y_j) = q_j = \sum_i r_{ij}$.
Tinjau selisih $H(X, Y) - H(X) - H(Y)$:
\begin{align*}
H(X, Y) - H(X) - H(Y) &= -\sum_{i}\sum_j r_{ij}\log_2 r_{ij} + \sum_i p_i \log_2 p_i + \sum_j q_j \log_2 q_j \\
&= -\sum_i \sum_j r_{ij}\log_2 r_{ij} + \sum_i \left(\sum_j r_{ij}\right)\log_2 p_i + \sum_j \left(\sum_i r_{ij}\right)\log_2 q_j \\
&= \sum_i \sum_j r_{ij} \left( -\log_2 r_{ij} + \log_2 p_i + \log_2 q_j \right) \\
&= \sum_i \sum_j r_{ij} \log_2\left( \frac{p_i q_j}{r_{ij}} \right).
\end{align*}
Terapkan Ketaksamaan Jensen~\eqref{eq:jensen_app} dengan memilih bobot $a_{ij} = r_{ij}$ (yang memenuhi $\sum_{i,j} r_{ij} = 1$) dan titik $t_{ij} = \frac{p_i q_j}{r_{ij}}$:
\begin{align*}
\sum_i \sum_j r_{ij} \log_2\left( \frac{p_i q_j}{r_{ij}} \right) &\le \log_2\left( \sum_i \sum_j r_{ij} \cdot \frac{p_i q_j}{r_{ij}} \right) \\
&= \log_2\left( \sum_i \sum_j p_i q_j \right) = \log_2\left( \sum_i p_i \sum_j q_j \right) \\
&= \log_2(1 \times 1) = \log_2(1) = 0.
\end{align*}
Maka diperoleh:
\[
H(X, Y) - H(X) - H(Y) \le 0 \implies \mathbf{H(X, Y) \le H(X) + H(Y)}.
\]
\end{proof}

\noindent\textbf{Langkah 3: Kondisi Kesamaan (Bagian c)}
Berdasarkan kondisi kesamaan Ketaksamaan Jensen untuk fungsi cekung tegas, tanda kesamaan berlaku jika dan hanya jika seluruh nilai $t_{ij}$ sama (konstan):
\[
\frac{p_i q_j}{r_{ij}} = c \implies r_{ij} = \frac{1}{c} p_i q_j.
\]
Karena $\sum_{i,j} r_{ij} = 1$ dan $\sum_{i,j} p_i q_j = 1$, maka nilai konstanta haruslah $c = 1$.
Akibatnya:
\[
r_{ij} = p_i q_j \iff P(X = x_i, Y = y_j) = P(X = x_i) \cdot P(Y = y_j) \quad \text{untuk seluruh } (i, j).
\]
Ini adalah definisi bahwa variabel acak $X$ dan $Y$ \textbf{saling bebas (\textit{independen})}.

\noindent\textbf{Langkah 4: Pembuktian $H(X|Y) \le H(X)$ (Bagian d)}
Dari Aturan Rantai Entropi: $H(X, Y) = H(Y) + H(X|Y)$.
Substitusikan ke dalam ketaksamaan subaditivitas:
\[
H(Y) + H(X|Y) \le H(X) + H(Y) \implies \mathbf{H(X|Y) \le H(X)}.
\]
Artinya, mengetahui nilai variabel $Y$ secara rata-rata tidak pernah menambah ketidakpastian mengenai variabel $X$ (\textit{conditioning never increases entropy}).

\vspace{0.4cm}

% =========================================================================
\subsection{Soal 11: Jarak Unicity \& Kunci Palsu \texorpdfstring{($\bar{s}_n$)}{(s-n)} Sandi Hill \texorpdfstring{$3 \times 3$}{3x3}}

\begin{soalsulit}{Soal 11 (Tingkat Sulit) -- Grup Matriks $\mathrm{GL}(3, \mathbb{Z}_{26})$, Kunci Palsu, \& Jarak Unicity}
Diberikan sistem Sandi Hill dengan ukuran blok matriks $3 \times 3$ atas alfabet $\mathbb{Z}_{26}$. Teks sumber adalah bahasa Inggris dengan parameter $H_L = 1.25\text{ bit/huruf}$ ($R_L \approx 0.75$, $\log_2 26 \approx 4.7004\text{ bit}$).
\textbf{Pertanyaan}:
\begin{enumerate}[label=\alph*)]
    \item Hitunglah ukuran ruang kunci $|\mathcal{K}| = |\mathrm{GL}(3, \mathbb{Z}_{26})|$ secara eksak menggunakan Teorema Sisa Cina (\textit{Chinese Remainder Theorem}).
    \item Hitunglah nilai entropi kunci $H(K) = \log_2 |\mathcal{K}|$ dalam satuan bit.
    \item Hitunglah Jarak Unicity $n_0$ untuk Sandi Hill $3 \times 3$ tersebut (dalam satuan jumlah karakter dan dalam satuan jumlah blok trigram).
    \item Misalkan seorang kriptanalis mencegat teks sandi sepanjang $n = 6$ karakter (2 blok trigram). Hitunglah batas bawah rata-rata banyaknya \textbf{kunci palsu ($\bar{s}_6$)} yang masih tersisa. Bandingkan hasilnya jika panjang teks sandi yang dicegat adalah $n = 15$ karakter (5 blok trigram).
\end{enumerate}
\end{soalsulit}

\noindent\textbf{Pembahasan Lengkap}:

\noindent\textbf{Langkah 1: Menghitung Ruang Kunci $|\mathrm{GL}(3, \mathbb{Z}_{26})|$ (Bagian a)}
Karena $26 = 2 \times 13$ dengan $\gcd(2, 13) = 1$, berdasarkan Teorema Sisa Cina:
\[
|\mathrm{GL}(3, \mathbb{Z}_{26})| = |\mathrm{GL}(3, \mathbb{Z}_2)| \times |\mathrm{GL}(3, \mathbb{Z}_{13})|.
\]
Rumus umum kardinalitas grup linear umum atas lapangan hingga $\mathbb{F}_q$:
\begin{equation}
|\mathrm{GL}(k, \mathbb{F}_q)| = (q^k - 1)(q^k - q)(q^k - q^2).
\end{equation}
\begin{enumerate}
    \item Untuk $q = 2, k = 3$:
    \[
    |\mathrm{GL}(3, \mathbb{Z}_2)| = (2^3 - 1)(2^3 - 2)(2^3 - 4) = (7)(6)(4) = \mathbf{168}.
    \]
    \item Untuk $q = 13, k = 3$:
    \begin{align*}
    |\mathrm{GL}(3, \mathbb{Z}_{13})| &= (13^3 - 1)(13^3 - 13)(13^3 - 13^2) \\
    &= (2197 - 1)(2197 - 13)(2197 - 169) \\
    &= (2196)(2184)(2028).
    \end{align*}
    Hitung nilai perkaliannya:
    \[
    2196 \times 2184 = 4.796.064 \implies 4.796.064 \times 2028 = \mathbf{9.726.417.792}.
    \]
    \item Total ukuran ruang kunci:
    \[
    |\mathcal{K}| = 168 \times 9.726.417.792 = \mathbf{1.634.038.189.056} \approx \mathbf{1.634 \times 10^{12}}.
    \]
\end{enumerate}

\noindent\textbf{Langkah 2: Menghitung Entropi Kunci $H(K)$ (Bagian b)}
Kunci diasumsikan dipilih seragam:
\begin{align*}
H(K) &= \log_2(1.634038 \times 10^{12}) = \log_2(168) + \log_2(9.726.417.792) \\
&= \frac{\ln 168}{\ln 2} + \frac{\ln 9726417792}{\ln 2} \approx 7.3923 + 33.1789 = \mathbf{40.5712\text{ bit}}.
\end{align*}

\noindent\textbf{Langkah 3: Menghitung Jarak Unicity $n_0$ (Bagian c)}
Penyebut rumus unicity:
\[
R_L \log_2 |\mathcal{P}| \approx 0.75 \times 4.7004 = 3.5253\text{ bit/karakter}.
\]
Maka jarak unicity:
\[
n_0 \approx \frac{H(K)}{R_L \log_2 |\mathcal{P}|} \approx \frac{40.5712}{3.5253} \approx 11.5086 \approx \mathbf{12\text{ karakter}}.
\]
Dalam ukuran blok: Karena 1 blok berisi 3 karakter, maka $n_0$ setara dengan:
\[
\left\lceil \frac{11.5086}{3} \right\rceil = \mathbf{4\text{ blok trigram}}.
\]

\noindent\textbf{Langkah 4: Perhitungan Kunci Palsu $\bar{s}_n$ (Bagian d)}
Gunakan rumus batas bawah ekspektasi kunci palsu Shannon (Teorema 9):
\begin{equation}
\bar{s}_n \ge \frac{|\mathcal{K}|}{|\mathcal{P}|^{n R_L}} - 1 = 2^{\log_2 |\mathcal{K}| - n R_L \log_2 26} - 1 = 2^{40.5712 - 3.5253 n} - 1.
\end{equation}
\begin{itemize}
    \item \textbf{Untuk $n = 6$ karakter (2 blok)}:
    \[
    \text{Eksponen} = 40.5712 - (3.5253 \times 6) = 40.5712 - 21.1518 = 19.4194.
    \]
    \[
    \bar{s}_6 \ge 2^{19.4194} - 1 \approx 703.958 - 1 \approx \mathbf{703.957\text{ kunci palsu}}.
    \]
    \textit{Makna}: Teks sandi sepanjang 6 huruf masih menyisakan lebih dari 700 ribu kunci palsu yang menghasilkan teks terbaca. Kunci belum dapat dipastikan secara unik.
    \item \textbf{Untuk $n = 15$ karakter (5 blok)}:
    \[
    \text{Eksponen} = 40.5712 - (3.5253 \times 15) = 40.5712 - 52.8795 = -12.3083.
    \]
    \[
    \bar{s}_{15} \ge 2^{-12.3083} - 1 \approx 0.000199 - 1 \implies \mathbf{\bar{s}_{15} \approx 0\text{ kunci palsu}}.
    \]
    \textit{Makna}: Karena $n = 15 > n_0 \approx 12$, seluruh kunci palsu telah tereleminasi dan kunci asli teridentifikasi secara tunggal.
\end{itemize}

\vspace{0.4cm}

% =========================================================================
\subsection{Soal 12: Analisis SPN Non-Idempoten, Confusion-Diffusion, \& Shannon}

\begin{soalsulit}{Soal 12 (Tingkat Sulit) -- Analisis Sistem Idempoten vs SPN Non-Idempoten Shannon}
Dalam makalah monumentalnya (1949), Claude Shannon meletakkan dua prinsip utama dalam perancangan cipher blok praktis yang aman: \textit{Confusion} dan \textit{Diffusion}, serta mengusulkan sistem perkalian tak-idempoten (\textit{non-idempotent product ciphers}).
\textbf{Pertanyaan}:
\begin{enumerate}[label=\alph*)]
    \item Tinjau iterasi perkalian dari sandi substitusi monoalfabetik $S^k = \underbrace{S \times S \times \dots \times S}_{k \text{ kali}}$ atas $\mathbb{Z}_{26}$. Tunjukkan bahwa ruang kunci efektif dari $S^k$ tidak bertambah dan tetap bernilai $26!$, serta jelaskan mengapa cipher idempoten rentan terhadap kriptoanalisis frekuensi.
    \item Jelaskan secara matematis definisi prinsip \textbf{\textit{Confusion}} dan \textbf{\textit{Diffusion}} menurut Shannon, serta sebutkan komponen modern yang merepresentasikannya pada arsitektur \textit{Substitution-Permutation Network} (SPN).
    \item Buktikan secara logis mengapa satu ronde fungsi SPN ($R(x) = P(S(x \oplus K))$) yang dikomposisikan secara berulang ($R^k$) \textbf{TIDAK bersifat idempoten} ($S_{\text{SPN}}^2 \neq S_{\text{SPN}}$).
    \item Mengapa sifat non-idempotensi pada SPN multi-ronde mampu menggagalkan serangan diferensial (\textit{Differential Cryptanalysis}) dan serangan linear (\textit{Linear Cryptanalysis}) pada cipher modern seperti AES?
\end{enumerate}
\end{soalsulit}

\noindent\textbf{Pembahasan Lengkap}:

\noindent\textbf{Langkah 1: Analisis Kunci Efektif $S^k$ \& Kerentanan Frekuensi (Bagian a)}
Misalkan $S_i$ adalah enkripsi permutasi dengan kunci $\pi_i \in S_{26}$. Enkripsi gabungan berantai adalah komposisi:
\[
e_{(\pi_1, \dots, \pi_k)}(x) = (\pi_k \circ \pi_{k-1} \circ \dots \circ \pi_1)(x).
\]
Karena himpunan seluruh permutasi $S_{26}$ membentuk grup berhingga di bawah operasi komposisi fungsi ($\circ$), maka hasil komposisi $\pi_{\text{total}} = \pi_k \circ \dots \circ \pi_1$ pasti merupakan suatu elemen permutasi tunggal di dalam $S_{26}$.
\begin{itemize}
    \item Meskipun penyerang melihat ada $k$ buah kunci $(\pi_1, \dots, \pi_k)$ dengan ukuran kombinasi $(26!)^k$, ruang kunci efektifnya tetap berukuran $|S_{26}| = 26! \approx 4.03 \times 10^{26}$.
    \item \textbf{Kerentanan Frekuensi}: Transformasi akhir $\pi_{\text{total}}$ tetap merupakan pemetaan satu-ke-satu sederhana pada alfabet ($x \mapsto \pi_{\text{total}}(x)$). Akibatnya, histogram distribusi frekuensi kemunculan karakter pada ciphertext tetap identik persis dengan histogram plaintext (hanya label hurufnya yang tertukar). Kriptanalis dapat memecahkan sandi $S^k$ dalam sekejap hanya dengan mencocokkan frekuensi huruf tanpa memedulikan nilai $k$.
\end{itemize}

\noindent\textbf{Langkah 2: Definisi Matematis Confusion \& Diffusion Shannon (Bagian b)}
\begin{enumerate}
    \item \textbf{\textit{Confusion} (Kebingungan)}:
    \begin{itemize}
        \item \textit{Definisi}: Membuat hubungan statistik antara nilai kunci rahasia $K$ dan ciphertext $C$ menjadi sekompleks dan senon-linear mungkin.
        \item \textit{Komponen SPN}: Direalisasikan oleh \textbf{S-Box (Kotak Substitusi non-linear)}. S-Box memetakan input $m$-bit ke output $n$-bit melalui fungsi Boolean berderajat non-linear tinggi sehingga tidak dapat didekati dengan persamaan aljabar linear.
    \end{itemize}
    \item \textbf{\textit{Diffusion} (Penyebaran)}:
    \begin{itemize}
        \item \textit{Definisi}: Menyebarkan struktur statistik dan redundansi dari plaintext $P$ ke seluruh bagian ciphertext $C$. Jika 1 bit teks asal diubah, rata-rata separuh bit teks sandi harus berubah secara acak (\textit{Strict Avalanche Criterion}).
        \item \textit{Komponen SPN}: Direalisasikan oleh \textbf{P-Box (Lapisan Permutasi / Difusi Linear)} atau perkalian matriks biner MDS (seperti \textit{MixColumns} pada AES).
    \end{itemize}
\end{enumerate}

\noindent\textbf{Langkah 3: Pembuktian SPN Tidak Idempoten (Bagian c)}
Tinjau operasi satu ronde SPN:
\[
R_1(x) = P(S(x \oplus K_1)).
\]
Komposisi dua ronde dengan dua subkunci independen $K_1, K_2$:
\[
R_2(R_1(x)) = P\Big( S\big( P(S(x \oplus K_1)) \oplus K_2 \big) \Big).
\]
\begin{proof}
\begin{enumerate}
    \item S-Box $S$ adalah pemetaan non-linear (derajat aljabar $d \ge 2$).
    \item Komposisi $S \circ P \circ S$ melibatkan evaluasi fungsi non-linear atas argumen yang merupakan keluaran fungsi non-linear sebelumnya yang telah diacak posisinya oleh permutasi $P$.
    \item Derajat polinomial aljabar dari fungsi komposisi dua ronde melonjak menjadi $d_{\text{total}} = d^2$.
    \item Bentuk $P(S(P(S(x \oplus K_1)) \oplus K_2))$ secara matematis tidak dapat difaktorkan atau disederhanakan menjadi bentuk satu ronde tunggal $P(S(x \oplus K'))$ untuk sembarang pemilihan kunci $K'$.
\end{enumerate}
Oleh karena itu, transformasi dua ronde tidak berada dalam keluarga transformasi satu ronde, sehingga terbukti bahwa:
\[
\mathbf{S_{\text{SPN}}^2 \neq S_{\text{SPN}} \quad (\text{TIDAK Idempoten})}.
\]
\end{proof}

\noindent\textbf{Langkah 4: Ketahanan Terhadap Kriptoanalisis Diferensial \& Linear (Bagian d)}
\begin{itemize}
    \item \textbf{Kriptoanalisis Diferensial}: Mengeksploitasi propagasi selisih input $\Delta x$ menuju selisih output $\Delta y$ dengan probabilitas $\mathrm{DP} = P(S(x \oplus \Delta x) \oplus S(x) = \Delta y)$. Pada 1 ronde, terdapat diferensial lokal dengan probabilitas tinggi. Namun karena sistem tidak idempoten dan memiliki difusi $P$, selisih bit menyebar ke banyak S-Box aktif di ronde berikutnya (\textit{active S-boxes}). Probabilitas karakteristik diferensial gabungan meluruh secara eksponensial:
    \[
    p_{\text{diferensial}} = \prod_{r=1}^R p_r \ll 2^{-128}.
    \]
    \item \textbf{Kriptoanalisis Linear}: Mengeksploitasi aproksimasi linear $a \cdot x \oplus b \cdot y = c \cdot K$ dengan bias $\epsilon$. Berdasarkan \textit{Piling-Up Lemma} Matsui, pengulangan $R$ ronde tak-idempoten dengan S-box aktif menyebabkan bias gabungan meluruh menjadi:
    \[
    \epsilon_{\text{total}} = 2^{m-1} \prod_{i=1}^m \epsilon_i \to 0.
    \]
    Ketika $\epsilon_{\text{total}} < 2^{-128}$, data yang dibutuhkan untuk menyerang melebihi seluruh ruang sampel plaintext yang ada ($2^{128}$), sehingga cipher terbukti aman secara komputasi.
\end{itemize}

\vspace{0.6cm}

% =========================================================================
\section{Tabel Rangkuman Parameter \& Kunci Jawaban Singkat}
% =========================================================================

\begin{table}[H]
\centering
\small
\renewcommand{\arraystretch}{1.5}
\begin{tabular}{c c p{4.2cm} p{8.8cm}}
\toprule
\textbf{No} & \textbf{Tingkat} & \textbf{Topik Materi} & \textbf{Hasil Akhir / Jawaban Kunci} \\
\midrule
1 & Mudah & Entropi Shannon & $H(X) = 1.75\text{ bit}, \; H_{\max} = 2.0\text{ bit}, \; \eta = 87.5\%, \; R = 12.5\%$. \\
2 & Mudah & Pengkodean Huffman & $l(f) = 1.90\text{ bit/simbol}, \; H(X) = 1.8465\text{ bit} \implies H \le l(f) < H+1$. \\
3 & Mudah & Jarak Unicity Sederhana & $R_L = 72.34\%, \; n_0(\text{Geser}) \approx 2\text{ huruf}, \; n_0(\text{Affine}) \approx 3\text{ huruf}$. \\
4 & Mudah & Idempotensi Sandi Geser & $S^2 = S$ terbukti; ruang kunci efektif tetap 26 (tidak menambah keamanan). \\
\midrule
5 & Sedang & Uji Keamanan Sempurna & $p(c_j) = 1/3, \; p(m_i|c_j) = p(m_i) \implies \text{Memenuhi Perfect Secrecy}$. \\
6 & Sedang & Key Equivocation & $H(P) = 1.5, \; H(K) = H(C) = 1.585, \; H(K|C) = H(P|C) = 1.5\text{ bit}$. \\
7 & Sedang & Jarak Unicity Vigenère / Permutasi & $n_0(\text{Vigenère}, m=8) \approx 11\text{ huruf}; \; n_0(\text{Permutasi}, S_8) \approx 5\text{ huruf}$. \\
8 & Sedang & Komposisi Sandi Affine & $M \times S = S \times M = \text{Affine}$; $\text{Affine}^2 = \text{Affine}$; Sub $\times$ Trans tak-komutatif. \\
\midrule
9 & Sulit & Syarat Shannon $|\mathcal{K}| \ge |\mathcal{P}|$ & Terbukti via Bayes $p(x|y)>0 \implies \exists K$, injeksi $\mathcal{C} \hookrightarrow \mathcal{K}$ dan $\mathcal{P} \hookrightarrow \mathcal{C}$. \\
10 & Sulit & Subaditivitas via Jensen & $f''(t) = -\frac{1}{t^2\ln 2} < 0$; Jensen $\implies H(X,Y) \le H(X)+H(Y) \iff X \perp Y$. \\
11 & Sulit & Sandi Hill $3 \times 3$ ($\mathrm{GL}$) & $|\mathcal{K}| \approx 1.634 \times 10^{12}, \; H(K) \approx 40.57\text{ bit}, \; n_0 \approx 12\text{ huruf}, \; \bar{s}_6 \approx 703.957$. \\
12 & Sulit & SPN Non-Idempoten & $S^k$ tereduksi ke $S_{26}$; SPN $R^2 \neq R$ (derajat non-linear melonjak $d^R$). \\
\bottomrule
\end{tabular}
\end{table}

\end{document}
